\documentclass{situationist-preprint}

\title{Posthumous Harm}
\msid{EXI-13-b}
\msfamily{EXI-13}
\msdiscipline{Meaning, Death and the Absurd}
\mstype{Companion manuscript}
\mskeywords{posthumous harm; duties to the dead; reliance; entitlement; surviving interests}

\begin{document}
\maketitle

\begin{abstract}
EXI-13-a grounds duties to the dead in two principles: surviving interests, whose thwarting is bad for the ante-mortem person, and reliance, which makes the community's practice of honoring the wishes of the dead a source of entitlements. We separate the two. Posthumous harm is a setback to a surviving interest; posthumous wrong is a breach of an obligation owed to an entitlement. The distinction answers the question, left open by ETH-16-b, whether the dead have entitlements: the dead as such have none, but the entitlements a person acquired while alive are not extinguished by her death. It also explains why there can be harm without wrong and wrong without much harm, and why both weaken over time.
\end{abstract}

\section{Harm and wrong}

Following \citet{Feinberg1984} and \citet{Pitcher1984}, EXI-13-a takes the subject of posthumous misfortune to be the ante-mortem person, whose interests in what happens after her death are thwarted or fulfilled by posthumous events (Lemma~3.1 of EXI-13-a). Its Principle~P1 says when such an event is bad for her; its Principle~P2 and Assumption~2.2 say when the living owe her something. These are different conditions.

\begin{definition}\label{d:hw}
A posthumous event \emph{harms} a person if it thwarts an interest of hers whose object lies after her death. A posthumous act \emph{wrongs} her if it breaches an obligation owed to an entitlement she acquired while alive.
\end{definition}

\begin{proposition}\label{p:apart}
Posthumous harm and posthumous wrong come apart in both directions.
\end{proposition}

\begin{proof}
Harm without wrong: a scientist's surviving interest in the confirmation of her theory is thwarted if, after her death, the theory is refuted by honest work. By Principle~P1 of EXI-13-a this is bad for her; but no one was obligated not to refute it, so no entitlement is breached. Wrong without much harm: a person reasonably wished that her letters be burnt and cared little about it; a community that gave her reliance on its honoring such wishes (Assumption~2.2 of EXI-13-a) wrongs her by publishing them, though the interest thwarted is slight. By Assumption~2.3 of EXI-13-a, the obligation then has slight weight against competing reasons, but it is not nothing.
\end{proof}

\section{Do the dead have entitlements?}

ETH-16-b argues that a deathbed promise creates an obligation by Principle~P1 of ETH-16-a when it is made, and leaves to EXI-13 the question whether the obligation persists after the promisee's death, that is, whether the dead can have entitlements.

\begin{proposition}\label{p:entitle}
The dead, considered as what remains after death, have no entitlements. The entitlements that a person acquired while alive are not extinguished by her death, and obligations owed to them persist after it.
\end{proposition}

\begin{proof}
An entitlement is created by an act of the obligor directed at a person who exists when it is created (Principle~P2 of EXI-13-a). Nothing is created in favor of a corpse, and after death no new reliance can be induced; so the dead as such acquire no entitlements. But by Principle~P2 of EXI-13-a, an obligation grounded in reliance persists until it is discharged, released or defeated. Death is none of these: it does not perform what was promised, it is not a release by the promisee, and it is not a weightier reason. So the obligation persists, and it is owed, by Lemma~3.1 of EXI-13-a, to the ante-mortem person who held the entitlement.
\end{proof}

This answers ETH-16-b in the sense it requires: the deathbed promise binds after the promisee's death, and it binds as a promise, since the entitlement on which it rests was created by the promise.

\begin{objection}[No subject at the time]
An obligation is owed to someone. After her death there is no one to whom it is owed, and no one who can be harmed \citep{Taylor2005}.
\end{objection}

\begin{reply}
The objection assumes that an obligation, or a harm, must have a subject who exists when it is breached or suffered. EXI-02 rejects this for harm: deprivation does not require a subject at the time of loss, and on the tenseless view the posthumous event is simply a fact about how her life went. The same holds for obligation: a promise to someone now unreachable, or asleep, or unaware, still binds, and what is required is that the promisee existed when the entitlement was created.
\end{reply}

\section{Weakening}

Both conditions weaken over time, and for the same reason. The interests whose setback is harm are, by Lemma~3.3 of EXI-13-a, progressively completed, extinguished or overtaken; and by Assumption~2.3 of EXI-13-a, the weight of an obligation owed to an entitlement is proportional to the interests it protects. So as harm diminishes, the weight of the wrong diminishes with it. The duty of Theorem~3.4 of EXI-13-a does not lapse, but posthumous harm and posthumous wrong fade together.

\begin{thebibliography}{9}\small
\bibitem[Feinberg(1984)]{Feinberg1984} Feinberg, J. (1984). \emph{Harm to Others}. Oxford University Press.
\bibitem[Pitcher(1984)]{Pitcher1984} Pitcher, G. (1984). The misfortunes of the dead. \emph{American Philosophical Quarterly}, 21(2), 183--188.
\bibitem[Taylor(2005)]{Taylor2005} Taylor, J.~S. (2005). The myth of posthumous harm. \emph{American Philosophical Quarterly}, 42(4), 311--322.
\end{thebibliography}

\end{document}
